% Generated draft: requires mathematical and source review.
\documentclass[11pt]{article}
\usepackage[margin=25mm]{geometry}
\usepackage{fontspec}
\setmainfont{Latin Modern Roman}
\usepackage{amsmath,amssymb,mathtools,unicode-math}
\setmathfont{Latin Modern Math}
\usepackage{enumitem,needspace,xurl,hyperref}
\setlength{\parindent}{0pt}
\setlength{\parskip}{5pt}
\setlength{\emergencystretch}{3em}
\newcommand{\N}{\mathbb{N}}
\newcommand{\Z}{\mathbb{Z}}
\newcommand{\Q}{\mathbb{Q}}
\newcommand{\R}{\mathbb{R}}
\newcommand{\C}{\mathbb{C}}
\newcommand{\F}{\mathbb{F}}
\newcommand{\A}{\mathbb{A}}
\newcommand{\PP}{\mathbb{P}}
\newcommand{\E}{\mathbb{E}}
\newcommand{\eps}{\varepsilon}
\newcommand{\dd}{\,\mathrm{d}}
\begin{document}
\section*{UnsolvedMath ID: 700383}
\textbf{Source problem code:} catalogue entry
\subsection{Definitions and mathematical statement}
Let \(k\) be a commutative field, let \(n\geq 1\) be an integer, and let \(L=\operatorname{GL}_n(k)\), the group of invertible \(n\times n\) matrices over \(k\). These are the ambient hypotheses explicitly supplied in part (a); their use in part (b) is a contextual interpretation. Say that a subgroup \(H\leq L\) meets every conjugacy class if
\[
\forall x\in L\quad \exists g\in L\quad gxg^{-1}\in H.
\]
A flag in \(k^n\) is a chain
\[
0=V_0<V_1<\cdots<V_r=k^n
\]
of \(k\)-linear subspaces. Under the usual flag-stabilizer interpretation, a parabolic subgroup of \(L\) is the full stabilizer
\[
\operatorname{Stab}_L(V_\bullet)=\{g\in L:gV_i=V_i\text{ for every }i\}.
\]
This convention includes \(L\) itself, using the flag \(0<k^n\). The source does not specify its convention for parabolic subgroups over arbitrary fields.

\textbf{Part (b), the remaining question in the supplied record.} Is it true, for every commutative field \(k\), every integer \(n\geq1\), and every subgroup \(H\leq\operatorname{GL}_n(k)\), that
\[
\bigl(\forall x\in\operatorname{GL}_n(k)\ \exists g\in\operatorname{GL}_n(k):gxg^{-1}\in H\bigr)
\ \Longrightarrow\ 
\bigl(\exists V_\bullet:H=\operatorname{Stab}_{\operatorname{GL}_n(k)}(V_\bullet)\bigr)?
\]
The source also asks how far the corresponding assertion holds for subgroups of other groups of Lie type. It does not specify the groups, fields, or interpretation of parabolic subgroup for this extension, so that request does not determine a single quantified statement.

\textbf{Part (a), retained as archived context.} A subgroup \(G\leq\operatorname{GL}_n(k)\) is irreducible if its only invariant \(k\)-linear subspaces of \(k^n\) are \(0\) and \(k^n\). The original classification question asks for all irreducible subgroups \(G\) meeting every conjugacy class of \(\operatorname{GL}_n(k)\). Neumann conjectured that every such \(G\) satisfies
\[
G=\operatorname{GL}_n(k)
\quad\text{or}\quad
\left\{
\begin{gathered}
n=2,\qquad \operatorname{char}(k)=2,\\
k\text{ is quadratically closed},\\
G\text{ is conjugate in }\operatorname{GL}_2(k)\text{ to }M(k),
\end{gathered}
\right.
\]
where
\[
M(k)=
\left\{\begin{pmatrix}\alpha&0\\0&\beta\end{pmatrix}:\alpha,\beta\in k^\times\right\}
\ \cup\ 
\left\{\begin{pmatrix}0&\alpha\\\beta&0\end{pmatrix}:\alpha,\beta\in k^\times\right\},
\qquad k^\times=k\setminus\{0\}.
\]
Here quadratically closed is understood to mean that every polynomial of degree two over \(k\) has a root in \(k\), hence splits over \(k\). The supplied archive reports that the conjecture in part (a) was refuted by S. A. Zyubin. This is a supplied historical assessment, not an assertion that the conjecture holds or that the original classification question has been completely answered.
\subsection{Short English statement}
Part (b) asks whether every subgroup of \(\operatorname{GL}_n(k)\) meeting every conjugacy class is a full flag stabilizer, and asks for analogous statements in other groups of Lie type. The field hypotheses are inherited contextually from part (a), and the scope of the Lie type extension is unspecified. The archived part (a) asks for a classification of irreducible such subgroups and records Neumann's conjecture, which the supplied source reports as refuted.
\subsection{Sources}
\begin{itemize}
\item Local UnsolvedMath registry, permanent ID \(700383\), title \emph{Kourovka Notebook Problem 6.38}, exact source problem code \texttt{catalogue entry}; supplied upstream revision \texttt{d371ee917ac9647fbb06cd6fc529deec6f3dda82}.
\item P. M. Neumann, Problem 6.38, \emph{The Kourovka Notebook (Unsolved Problems in Group Theory)}, 6th issue (1978), as represented in the supplied record, updated in the 21st issue (2026), reference \texttt{arXiv:1401.0300v48}. The supplied collection description names E. I. Khukhro and V. D. Mazurov as editors.
\item The supplied Archive of Solved Problems entry attributes the refutation of the conjecture in part (a) to S. A. Zyubin, \emph{Algebra and Logic} \textbf{45}, no. 5 (2006), 296--305. No article title is supplied.
\item The supplied status review is dated \texttt{2026-10-16}, labels the record \texttt{partially\_solved}, and says that part (a) appears in the solved archive while remaining parts appear in the main section. This assessment is supplied dated curation, not independently verified source or literature checking.
\end{itemize}
\end{document}
